%

\subsection{A Function-Space Perspective on Variational Inference in Bayesian Neural Networks}
%
\label{sec:theory}

Instead of seeking to infer an approximate posterior distribution over parameters, we frame variational inference in stochastic neural networks as inferring an approximation to the posterior distribution over \emph{functions} $ p_{f(\cdot \,; \bTheta) | \calD}$ induced by the posterior distribution over parameters $p_{\bTheta | \calD}$, that is,
\begin{align}
\SwapAboveDisplaySkip
    p_{f(\cdot \,; \bTheta) | \calD}(f(\cdot \,; \btheta) \vbar \calD)
    =
    \int_{\R^{P}} p_{\bTheta | \calD}(\btheta' \vbar \calD) \, \delta( f(\cdot \,; \btheta) - f(\cdot \,; \btheta') ) \dee \btheta' ,
\end{align}
where $\delta(\cdot)$ is the Dirac delta function~\citep{wolpert1993fsmap}.
Considering the prior distribution over functions $p_{f(\cdot \,; \bTheta)}$ induced by a prior distribution over parameters $p_{\bTheta}$,
\begin{align}
\SwapAboveDisplaySkip
    p_{f(\cdot \,; \bTheta)}(f(\cdot \,; \btheta))
    =
    \int_{\R^{P}} p_{\bTheta}(\btheta') \, \delta( f(\cdot \,; \btheta) - f(\cdot \,; \btheta') ) \dee \btheta' ,
\end{align}
and the variational distribution over functions $q_{f( \cdot \,; \bTheta)}$ induced by a variational distribution over parameters $q_{\bTheta}$,
\begin{align}
\SwapAboveDisplaySkip
    q_{f( \cdot \,; \bTheta)}(f(\cdot \,; \btheta))
    =
    \int_{\R^{P}} q_{\bTheta}({\btheta'}) \, \delta( f(\cdot \,; \btheta) - f(\cdot \,; \btheta') ) \dee \btheta' ,
\end{align}
we can express the problem of finding a posterior distribution over functions variationally as
\begin{align}
%
    \min_{q_{\bTheta} \in \calQ_{q_{\bTheta}}} \DD_{\textrm{KL}}( q_{f( \cdot \,; \bTheta)} \,\|\, p_{f(\cdot \,; \bTheta) | \calD} ) ,
\end{align}
which allows us to effectively incorporate meaningful prior information about the underlying data-generating process into training.
As discussed by~\citet{burt2021understanding}, this variational objective is guaranteed to be well-defined for suitably chosen prior distributions over functions.
Specifically, the \kld between two distributions over functions generated from different distributions over parameters applied to the same mapping (e.g., the same neural network architecture) is well-defined (i.e., finite) if the \kld between the distributions over parameters is finite, since, by the strong data processing inequality~\citep{Polyanskiy2017data},
\begin{align}
    \DD_\textrm{KL}(q_{f( \cdot \,; \bTheta)} \,\|\, p_{f(\cdot ; \bTheta)}) \leq \DD_\textrm{KL}(q_{\bTheta} \,\|\, p_{\bTheta}).
\end{align}
As a result, if $\DD_\textrm{KL}(q_{\bTheta} \,\|\, p_{\bTheta}) < \infty$, which is the case for finite-dimensional parameter vectors $\bTheta$ and $q_{\bTheta}$ absolutely continuous with respect to $p_{\bTheta}$, then the function-space \kld is finite and thus well-defined as a variational objective.
%

Hence, for a likelihood function defined on a finite set of training targets $\by_{\calD}$ and a suitably defined prior distribution over functions, we can express the variational problem above equivalently as the well-defined maximization problem \mbox{$\max_{q_{\bTheta} \in \calQ_{\btheta}} \mathcal{F}(q_{\bTheta})$} with
\begin{align}
\begin{split}
\label{eq:kld_function_space}
    %
    \mathcal{F}(q_{\bTheta})
    &
    \defines
    %
    \mathbb{E}_{q_{f(\bX_{\calD} ; \bTheta)}}[\log p_{\by | f(\bX ; \bTheta)}(\by_{\calD} \vbar f(\bX_{\calD} ; \btheta) )]
    %
    %
    - \DD_\textrm{KL}(q_{f(\cdot ; \bTheta)} \,\|\, p_{f(\cdot ; \bTheta)})
    %
    ,
\end{split}
\end{align}%
where $\DD_\textrm{KL}(q_{f(\cdot ; \bTheta)} \,\|\, p_{f(\cdot ; \bTheta)})$ is also a \kld between distributions over functions.

Unfortunately, evaluating the \kld in~\Cref{eq:kld_function_space} is in general intractable for arbitrary mappings $f$.
To obtain a tractable objective,~\citet{sun2019fbnn} showed that $\DD_\textrm{KL}(q_{f(\cdot ; \bTheta)} \,\|\, p_{f(\cdot ; \bTheta)})$ can be expressed as the supremum of the \kld from $q_{f(\cdot ; \bTheta)}$ to $p_{f(\cdot ; \bTheta)}$ over all \emph{finite} sets of evaluation points, resulting in the objective function
\begin{align}
\begin{split}
\label{eq:kld_fbnn}
    \mathcal{F}(q_{\bTheta})
    %
    =
    \mathbb{E}_{q_{f(\bX_{\calD} ; \bTheta)}}[\log p_{\by | f(\bX ; \bTheta)}(\by_{\calD} \vbar f(\bX_{\calD} ; \btheta ))]
    %
    %
    - \sup_{\bX \in \calX_{\mathbb{N}}} \hspace*{-3pt}\DD_\textrm{KL}(q_{f(\bX ; \bTheta)} \,\|\, p_{f(\bX ; \bTheta)} ) ,
    %
\end{split}
\end{align}%
where $\calX_{\mathbb{N}} \defines \bigcup_{n \in \mathbb{N}} \{ \bX \in \calX_{n} \vbar \calX_{n} \subseteq \mathbb{R}^{n \times D} \}$ is the collection of all finite sets of evaluation points.
However, this objective function is still challenging to optimize in practice:
The supremum cannot be obtained analytically and the \kld term itself is analytically intractable and difficult to estimate in high dimensions---even for a single evaluation point.
%
%
%

In the next section, we will describe an approximation and estimation procedure that allows scaling function-space variational inference to large neural networks and high-dimensional input data.

%
%
%
%
%



%

%
%
%
%
%
%
%
%
%
%
%
%
%

%
%
%

%
%
%
%
%

%
%
%
%
%

%
%
%
%

%



%
\section{Deriving a Tractable Function-Space Variational Objective}
\label{sec:fsvi_linearization}

The primary obstacle to computing the objective in~\Cref{eq:kld_function_space} is the \kld from $q_{f(\cdot ; \bTheta)}$ to $p_{f(\cdot ; \bTheta)}$.
There are two reasons why the \kld in~\Cref{eq:kld_fbnn} is intractable:
First, for \bnns or other non-linear models, we do not have access to the probability density functions of the multivariate distributions $q_{f(\bX ; \bTheta)}$ and $p_{f(\bX ; \bTheta)}$;
second, for all but extremely simple input spaces, we are unable to compute the supremum over all possible finite sets of evaluation points.
In the remainder of this section, we outline an approach for obtaining an estimator of a locally accurate approximation to the \kld that allows for scalable gradient-based optimization of~\Cref{eq:kld_fbnn}.

We first approach the problem of computing the \kld between two \bnns evaluated at a finite set of points.
To do so, we first derive tractable approximations to the distributions over functions $q_{f(\bX ; \bTheta)}$ and $p_{f(\bX ; \bTheta)}$
Next, we show that under these approximations, we are able to obtain a closed-form approximation to the \kld and describe a simple Monte Carlo estimator of the supremum in the function-space \kld.


\subsection{Approximating Distributions over Functions via Local Linearization}
\label{subsec:linearizations}

To obtain an approximation to the probability distributions of $q_{f(\bX ; \bTheta)}$ and $p_{f(\bX ; \bTheta)}$, we use a first-order Taylor expansion of the \emph{mapping} $f$ about the mean parameters of $q_{\bTheta}$ and $p_{\bTheta}$, respectively, and derive the induced distributions under the linearized mapping.

%
%
\label{prop:linearized_dist}
For a stochastic function $f(\cdot \,; \bTheta)$ defined in terms of stochastic parameters \mbox{$\bTheta$} distributed according to distribution $g_{\bTheta}$ with $\mathbf{m} \defines \E_{g_{\bTheta}}[\bTheta]$ and $\mathbf{S} \defines \text{Cov}_{g_{\bTheta}}[\bTheta]$, we denote the linearization of the stochastic function $f(\cdot \,; \bTheta)$ about $\mathbf{m}$ by
\begin{align}
%
\label{eq:linearization}
    f(\cdot \,; \bTheta) \approx \flin(\cdot \,; \mathbf{m}, \bTheta) \defines f(\cdot \,; \mathbf{m}) + \jac(\cdot \,; \mathbf{m})(\bTheta - \mathbf{m}),
\end{align}
where $\jac(\cdot \,; \mathbf{m}) \defines (\partial f(\cdot \,; \bTheta) / \partial \bTheta)|_{\bTheta = \mathbf{m}}$ is the Jacobian of $f(\cdot \,; \bTheta)$ evaluated at $\bTheta = \mathbf{m}$, and the mean and covariance of the distribution over the linearized mapping $\flin$ at \mbox{$\bX, \bX' \in \calX$} are given by
\begin{align}
%
    \E[\flin(\bX ; \bTheta)]
    &
    =
    f(\bX; \mathbf{m})
    \\
    %
    \textrm{{Cov}}[\flin(\bX ; \bTheta), \flin(\bX' ; \bTheta)]
    &
    =
    \jac(\bX ; \mathbf{m}) \mathbf{S} \jac(\bX', \mathbf{m})^\top .
\end{align}
For a derivation of this result, see~\Cref{appsec:proofs}.
Since Gaussianity is preserved under affine transformations, if $g_{\bTheta}$ is a multivariate Gaussian distribution with mean $\mathbf{m}$ and diagonal co-variance $\mathbf{S}$, then the distribution $\tilde{g}$ over $\flin(\bX \,; \bTheta)$ is given by
%
\begin{align}
%
\label{eq:dist_lin}
    \tilde{g}_{\flin(\bX ; \mathbf{m}, \bTheta)} = \mathcal{N}(f(\bX ; \mathbf{m}), \jac(\bX ; \mathbf{m}) \mathbf{S} \jac(\bX ; \mathbf{m})^\top) .
\end{align}
For stochastic functions parameterized by many millions of parameters, obtaining the covariance of $\tilde{g}_{\flin(\bX ; \bTheta)}$---which requires computing an inner product of two Jacobian matrices---can be computationally expensive.
Instead of computing the distribution over the linearized mapping exactly, we can construct a suitable Monte Carlo estimator.
To do so, we consider a partition of the set of parameters into sets $\alpha$ and $\beta$ (with $|\beta| \ll |\alpha|$) and note that the linearized mapping can then be expressed as
\begin{align}
\label{eq:linearization_decomposition}
    \flin(\cdot \,; \mathbf{m}, \bTheta)
    =
    f(\cdot \,; \mathbf{m}) + \flin_{\alpha}(\cdot \,; \mathbf{m}, \bTheta_{\alpha}) + \jac_{\beta}(\cdot \,; \mathbf{m})(\bTheta_{\beta} - \mathbf{m}_{\beta}) ,
\end{align}
with
\begin{align}
\SwapAboveDisplaySkip
    \flin_{\alpha}(\cdot \,; \mathbf{m}, \bTheta_{\alpha}) \defines \jac_{\alpha}(\cdot \,; \mathbf{m})(\bTheta_{\alpha} - \mathbf{m}_{\alpha}) ,
\end{align}
where $\jac_{\alpha}(\cdot \,; \mathbf{m})$ and $\jac_{\beta}(\cdot \,; \mathbf{m})$ are the columns of the Jacobian matrix corresponding to the sets of parameters $\alpha$ and $\beta$, respectively, and $\bTheta_{\alpha}$ and $\bTheta_{\beta}$ are the corresponding random parameter vectors.
Noting that~\Cref{eq:linearization_decomposition} expresses $\flin$ as a sum of (affine transformations of) random variables, we can use the fact that for independent Gaussian random variables $\bX$ and $\mathbf{Y}$, the distribution $h_{\mathbf{Z}}$ of $\mathbf{Z} = \bX + \mathbf{Y}$ is equal to the convolution of the distributions $h_{\bX}$ and $h_{\mathbf{Y}}$ to obtain an approximation to $\flin$.
In particular, we can show that if $g_{\bTheta}$ is a multivariate Gaussian distribution with \mbox{$\bTheta_{\alpha} \perp \bTheta_{\beta}$}, the distribution $\tilde{g}_{\flin(\bX ; \bTheta)}$ can be approximated by the Monte Carlo estimator
\begin{align}
\label{eq:dist_lin_approx}
    \hat{\tilde{g}}_{\flin(\bX ; \mathbf{m}, \bTheta)}
    =
    \frac{1}{R} \sum\nolimits_{j=1}^{R} \mathcal{N}\Big(f(\bX ; \mathbf{m}) + \flin_{\alpha}(\bX ; \mathbf{m}, \bTheta_{\alpha})^{(j)}, \jac_{\beta}(\bX ; \mathbf{m}) \mathbf{S}_{\beta} {\jac_{\beta}(\bX ; \mathbf{m})}^\top \Big) ,
\end{align}
where $g_{\bTheta_{\beta}} = \calN(\mathbf{m}_{\beta}, \mathbf{S}_{\beta})$ and samples $\flin_{\alpha}(\bX ; \mathbf{m}, \bTheta_{\alpha})^{(j)}$ are obtained by sampling parameters from the distribution \mbox{$g_{\bTheta_{\alpha}} =  \calN(\mathbf{m}_{\alpha}, \mathbf{S}_{\alpha})$}.
For a derivation of this result, see~\Cref{appsec:proofs}.
This estimator is biased for finite $K$ but converges to $\tilde{g}_{\flin(\bX ; \mathbf{m}, \bTheta)}$ as $R \rightarrow \infty$.
Similarly, for finite $R$, the smaller $ [\mathbf{S}_{\alpha}]_{ii} $, the more accurate and less biased the estimator will be.
In our empirical evaluation, we use a single Monte Carlo sample, $R=1$, to preserve Gaussianity and choose $\alpha$ to be the set of parameters in neural network layers \mbox{$1:L-1$} and $\beta$ to be the set of parameters in the final neural network layer.
%





\subsection{Approximating the Function-Space Kullback-Leibler Divergence}
\label{subsec:kl_approximation}

From~\Cref{subsec:linearizations}, we know that if $q_{\bTheta}$ and $p_{\bTheta}$ are both Gaussian distributions, then the induced distributions under the linearized mapping $\ftilde$ evaluated at a finite set of evaluation points will be Gaussian as well.
This means that for Gaussian variational and prior distributions over $\bTheta$, we can obtain locally accurate approximations to the induced distributions $q_{f(\cdot ; \bTheta)}$ to $p_{f(\cdot ; \bTheta)}$ and use them to approximate the \kld in the variational objective by $\DD_\textrm{KL}(\qtilde_{\ftilde(\bX ; \bTheta)} \,\|\, \ptilde_{\ftilde(\bX ; \bTheta)})$.
Moreover, for an isotropic Gaussian prior and a mean-field Gaussian variational distribution, $\DD_\textrm{KL}(\qtilde_{\ftilde(\bX ; \bTheta)} \,\|\, \ptilde_{\ftilde(\bX ; \bTheta)})$ is a \kld between two multivariate Gaussians and can be obtained analytically.

Using this approximation, we obtain an estimator of the variational objective given by
\begin{align}
\begin{split}
    %
    &\hspace*{-5pt}\tilde{\mathcal{F}}(q_{\bTheta})
    %
    \defines
    %
    \mathbb{E}_{q_{f(\bX_{\calD} ; \bTheta)}}[\log p_{\by | f(\bX ; \bTheta)}(\by_{\calD} \vbar f(\bX_{\calD} ; \btheta) )]
    %
    %
    - \sup_{\bX \in \calX_{\mathbb{N}}}\DD_\textrm{KL}(\qtilde_{\ftilde(\bX ; \bTheta)} \,\|\, \ptilde_{\ftilde(\bX ; \bTheta)})
    %
    ,
\end{split}
\end{align}%
where the arguments of the \kld have been replaced by the (locally accurate) approximations to the variational and prior distributions over functions evaluated at $\bX$, respectively.
Since the stochastic functions $\ftilde(\cdot \,; \bTheta)$ induced by $q_{\bTheta}$ and $p_{\bTheta}$ under the linearized mapping will be closer to the stochastic function under $f$ the smaller the variance of $q_{\bTheta}$ and $p_{\bTheta}$, respectively, the approximation to the \kld will be more accurate the smaller the variance of $q_{\bTheta}$ and $p_{\bTheta}$.

Next, we turn to computing the supremum.
Unlike~\citet{sun2019fbnn}, who consider the supremum as a separate optimization problem, we do not seek to compute the supremum by searching over points $\bX \in \calX_{\mathbb{N}}$ but instead propose to estimate the supremum at every gradient step via a simple finite-sample estimator.
Specifically, letting $I(\bX) \defines \DD_\textrm{KL}(\qtilde_{\ftilde(\bX ; \bTheta)} \,\|\, \ptilde_{\ftilde(\bX ; \bTheta)})$, we estimate $G = \sup_{\bX \in \calX_{\mathbb{N}}} I(\bX)$ using the Monte Carlo estimator
%
\begin{align}
\SwapAboveDisplaySkip
\label{eq:kl_estimator_max}
    \hat{G}(\calX_{\calC}^{S})
    =
    \max_{ \bX \in \calX_{\calC}^{S}} I(\bX) ,
    %
    %
\end{align}%
where \mbox{$\calX_{\calC}^{S} \defines \{ \cX^{(i)} \}_{i=1}^S$} is a collection of $S$ sets of \textit{context points} \mbox{$\cX^{(i)} \defines \{ \bx^{(j)} \}_{j=1}^{K}$} jointly sampled from a context distribution $p_{\calX_{\calC}}$.
Each context set $\cX^{(i)}$ can be viewed as a single Monte Carlo sample from the input space so that the estimator $\hat{G}(\calX_{\calC}^{S})$ provides an $S$-sample Monte Carlo estimate of the supremum.
%
While this estimator is crude and only provides a rough approximation to the true supremum, it encourages the variational distribution over functions to match the prior distribution over functions on the sets of context points.
%
The choice of the context distribution $p_{\calX_{\calC}}$ can be informed by knowledge about the prediction task and should be viewed as a problem-specific modeling choice.
Similarly, the numbers of samples $S$ and $K$ are hyperparameters to be optimized with a validation set.
For details on how $p_{\calX_{\calC}}$ is chosen for the empirical evaluation in~\Cref{sec:experiments}, see~\Cref{appsec:model_details}.
%


\subsection{Stochastic Estimation of the Approximate Function-Space Variational Objective}

Let $q_{\bTheta}$ be a Gaussian mean-field variational distribution,  let $p_{\bTheta}$ be an isotropic Gaussian prior, let $(\bX_{\calB}, \by_{\calB})$ be a mini-batch of the training data, and reparameterize $\bTheta$ as $\hat{\bTheta}(\bmu, \bSigma, \bepsilon^{(j)}) \defines \bmu + \bSigma \odot \bepsilon^{(j)}$.
Using the estimator $ \hat{G}(\calX_{\calC}^{S})$ defined above and estimating the expected log-likelihood via Monte Carlo sampling, we obtain a Monte Carlo estimator for the function-space variational objective:
\begin{align}
%
%
\label{eq:variational_objective_diag}
    %
    &
    \bar{\calF}(\bmu, \bSigma)
    %
    %
    %
    %
    =
    \frac{1}{M} \sum\nolimits_{j=1}^{M} \log p_{\by | f(\bX ; \bTheta)}(\by_{\calB} \vbar f(\bX_{\calB} ; \hat{\bTheta}(\bmu, \bSigma, \bepsilon^{(j)})) )
    %
    %
    - \max_{ \bX \in  \calX_{\calC}^{S} } {\DD_\textrm{KL}(\qtilde_{\ftilde(\bX ; \hat{\bTheta})} \,\|\, \ptilde_{\ftilde(\bX ; \hat{\bTheta})})}
    %
%
\end{align}
with $\bepsilon^{(j)} \sim \calN(\mathbf{0}, \mathbf{I}_{P})$ and \mbox{$\calX_{\calC}^{S}$} as defined above.
This Monte Carlo estimator is biased due to the linearization and context-set approximations but allows for scalable gradient-based stochastic optimization.
%
%

\textbf{Selection of Prior.}$~$
For all experiments that involve uncertainty quantification, we chose a prior distribution over parameters that induces a prior distribution over functions $p_{f(\cdot ; \bTheta)}$ and a prior predictive distribution that exhibits a high degree of predictive uncertainty at evaluation points from regions in input space where $p_{\calX_{\calC}}$ has non-zero support and, under smoothness constraints, on evaluation points in nearby regions.
For settings where prior information is encoded in data---for example, in the form of expert demonstrations of robotic manipulation tasks~\citep{rudner2021pathologies} or in the form of pre-trained networks in continual or transfer learning~\citep{Rudner2022sfsvi}---an empirical prior that reflects this information can be specified.
For further details, see~\Cref{appsec:model_details}.
%


\textbf{Selection of Context Distribution.}$~$
The distribution $p_{\calX_{\calC}}$ allows us to incorporate information about the data-generating process into training and encourage the variational distribution to match the prior over functions in relevant parts of the input space.
By taking advantage of the abundance of data available in real-world settings, context distributions can be constructed from large datasets like ImageNet~\citep{alexnet}, from small but diverse datasets like CIFAR-100, or by using any set of task-related unlabeled data.
In our experiments, we choose two types of context distributions.
One of the context distributions is constructed from the training data and only contains randomly sampled monochrome images, and one is constructed from a real-world dataset generated from a data distribution related to that of the training data.
For example, when training on FashionMNIST, we use KMNIST as the context distribution, and when training on CIFAR-10, we use CIFAR-100 as the context distribution.
For further details, see~\Cref{appsec:model_details}.
%
%

\textbf{Posterior Predictive Distribution.}$~$
After optimizing the variational objective with respect to the parameters of the variational distribution $q_{\bTheta}$, we use the fact that we can obtain function draws by sampling from the distribution over parameters to obtain an approximate posterior predictive distribution
\begin{align}
%
\begin{split}
    %
    \hspace*{-3pt}q(\by_{\ast} \vbar \bx_{\ast})
    &
    =
    \int p(\by_{\ast} \vbar f(\bx_{\ast} ; \btheta)) \, q_{f(\bx_{\ast} ; \bTheta)} \, \dee f(\bx_{\ast} ; \btheta)
    \\
    &
    \approx
    \frac{1}{M_{\ast}} \sum\nolimits_{j=1}^{M_{\ast}} p(\by_{\ast} \vbar f(\bx_{\ast} ; \bTheta^{(j)}))
    \quad
    \text{with}
    \quad
    \bTheta^{(j)} \sim q_{\bTheta} ,\hspace*{-5pt}
\end{split}
\end{align}%
where $M_{\ast}$ is the number of Monte Carlo samples used to estimate the predictive distribution.

%
%
%
%



%

%
%
%
%
%
%
%
%
%
%
%
%
%





%
%
%

%
%
%
%
%
%
%
%
%
%
%
%
%

%
%
%
%
%
%
%
%
%
%
%
%
%
%
%

%
%
%
%
%



%
